
\documentclass[11pt]{article}
\usepackage[letterpaper,margin=0.75in]{geometry}
\usepackage[T1]{fontenc}
\usepackage[utf8]{inputenc}
\usepackage{lmodern}
\usepackage{array}
\usepackage{longtable}
\usepackage{enumitem}
\usepackage{xcolor}
\usepackage{amssymb}
\usepackage{xurl}
\usepackage[hidelinks]{hyperref}
\setlength{\parindent}{0pt}
\setlength{\parskip}{6pt}
\emergencystretch=3em
\setlist[itemize]{leftmargin=1.2em,itemsep=2pt,topsep=2pt}
\setlist[enumerate]{leftmargin=1.4em,itemsep=2pt,topsep=2pt}
\renewcommand{\arraystretch}{1.25}
\newcommand{\checkbox}{$\Box$\hspace{0.4em}}
\newcommand{\ok}{\textbf{Success criteria:} }
\newcommand{\sourcebasis}{\section*{Official Source Basis}\begin{itemize}
\item Authentik Docker Compose installation: \url{https://docs.goauthentik.io/install-config/install/docker-compose/}
\item Authentik reverse proxy guidance: \url{https://docs.goauthentik.io/install-config/reverse-proxy/}
\item Authentik configuration reference: \url{https://docs.goauthentik.io/install-config/configuration/}
\item Authentik first steps: \url{https://docs.goauthentik.io/install-config/first-steps/}
\item Authentik backup and restore: \url{https://docs.goauthentik.io/sys-mgmt/ops/backup-restore/}
\item Authentik upgrade guidance: \url{https://docs.goauthentik.io/install-config/upgrade/}
\item Authentik monitoring: \url{https://docs.goauthentik.io/sys-mgmt/ops/monitoring/}
\item Authentik example flows: \url{https://docs.goauthentik.io/add-secure-apps/flows-stages/flow/examples/flows/}
\item Authentik service accounts: \url{https://docs.goauthentik.io/sys-mgmt/service-accounts/}
\item Authentik certificates: \url{https://docs.goauthentik.io/sys-mgmt/certificates/}
\end{itemize}}

\title{Authentik Security Hardening Checklist}\author{Prepared from official Authentik documentation}\date{May 6, 2026}\begin{document}\maketitle

\section*{Purpose}
This hardening checklist focuses on the controls most relevant to a Docker Compose Authentik deployment using the official documentation: secure proxying, trusted headers, MFA flows, Docker socket exposure, service-account least privilege, backups, upgrades, and safe troubleshooting.

\section*{Network and TLS Hardening}
\begin{itemize}
\item \checkbox Place Authentik behind an HTTPS-capable reverse proxy when exposing it through a public or shared network name.
\item \checkbox Pass \texttt{X-Forwarded-Proto}, \texttt{X-Forwarded-For}, \texttt{Host}, and WebSocket upgrade headers to Authentik.
\item \checkbox Use modern TLS configuration and disable SSL/TLS protocols older than TLS 1.3 where feasible.
\item \checkbox Configure \texttt{AUTHENTIK\_LISTEN\_\_TRUSTED\_PROXY\_CIDRS} when proxy headers are accepted from a proxy path outside the default trusted ranges.
\item \checkbox Use a certificate generated outside Authentik that matches the FQDN; prefer a publicly trusted certificate authority for proxy provider and brand use cases.
\end{itemize}

\section*{Secrets and Configuration Hardening}
\begin{itemize}
\item \checkbox Generate a strong \texttt{PG\_PASS} during initial deployment.
\item \checkbox Generate a strong \texttt{AUTHENTIK\_SECRET\_KEY} during initial deployment.
\item \checkbox Keep \texttt{.env} file access limited to administrators and automation accounts that require it.
\item \checkbox Use the worker dump-config command to verify applied configuration when needed.
\item \checkbox Avoid mounting \texttt{/etc/timezone} or \texttt{/etc/localtime} into Authentik containers; Authentik internal operations use UTC.
\end{itemize}

\section*{Docker Socket and Outpost Hardening}
\begin{itemize}
\item \checkbox Treat \texttt{/var/run/docker.sock} exposure to the worker container as privileged access.
\item \checkbox Use a Docker socket proxy when automatic outpost management is required and the organization wants reduced socket exposure.
\item \checkbox Remove the Docker socket mount and manually deploy/manage outposts when automatic outpost management is not required.
\item \checkbox Keep Authentik server and outpost versions matched during upgrades.
\end{itemize}

\section*{Identity and Access Hardening}
\begin{itemize}
\item \checkbox Use application/provider pairs deliberately; every protected application should have a clear provider configuration and access policy.
\item \checkbox Use bindings to restrict application visibility and access to intended users or groups.
\item \checkbox Use the two-factor login example flow or equivalent MFA design when additional verification is required.
\item \checkbox Force two-factor authentication by adjusting the Authenticator Validation Stage behavior where required.
\item \checkbox For recovery flows, prefer recovery designs that include MFA verification when the user population supports it.
\end{itemize}

\section*{Service Account Hardening}
\begin{itemize}
\item \checkbox Use service accounts only for machine-to-machine access and automation.
\item \checkbox Remember that service accounts cannot log in through the UI, cannot have a password, and use tokens exclusively.
\item \checkbox Prefer expiring tokens where operationally feasible.
\item \checkbox Revoke or rotate service account tokens when they are no longer needed.
\item \checkbox Apply least privilege: assign only the groups, permissions, applications, and resources required for the integration.
\end{itemize}

\section*{Backup, Upgrade, and Recovery Hardening}
\begin{itemize}
\item \checkbox Back up PostgreSQL regularly because it stores persistent Authentik data including users, policies, and configuration.
\item \checkbox Back up \texttt{/data} if file storage is local and not backed by S3.
\item \checkbox Back up \texttt{/certs} when certificates are relied on from the filesystem.
\item \checkbox Back up \texttt{/custom-templates} when UI customizations exist.
\item \checkbox Back up \texttt{/blueprints} when custom blueprints are used.
\item \checkbox Before upgrades, read release notes, back up PostgreSQL, and follow the documented upgrade sequence.
\item \checkbox Do not plan on downgrading Authentik; the official guidance states downgrades are unsupported.
\end{itemize}

\section*{Monitoring and Troubleshooting Hardening}
\begin{itemize}
\item \checkbox Monitor \texttt{/-/health/live/} for server liveness.
\item \checkbox Monitor \texttt{/-/health/ready/} for readiness and PostgreSQL connectivity.
\item \checkbox Monitor worker health with \texttt{ak healthcheck} inside the worker container.
\item \checkbox Monitor outposts with \texttt{/outpost.goauthentik.io/ping}.
\item \checkbox Scrape metrics on port 9300 only from intended monitoring networks; metrics are unauthenticated and non-exposed by default.
\item \checkbox Keep normal log level at \texttt{info} unless troubleshooting requires temporary escalation.
\item \checkbox Avoid \texttt{trace} logging unless absolutely necessary; trace logs can expose sensitive information including session cookies.
\end{itemize}

\sourcebasis

\end{document}
